%!TEX root = ../main.tex
\begin{center}

{\LARGE \bf Operational Optimization Methods Using Mathematical Models for a Power Grid with Photovoltaic Generation}
\\[0.7cm] 
{\Large YUMA OKUMURA}
\\[1.0cm]
{\LARGE \bf ABSTRACT}
\end{center}

In recent years, initiatives such as ``Carbon Neutrality by 2050'' and the ``GX2040 Vision'' have been promoted to address global warming, accelerating the deployment of renewable energy. While photovoltaics are anticipated to play a substantial role as distributed power sources, their fluctuations of weather-dependent output disrupt the grid's supply-demand balance. As a result, effective power management using battery storage systems has become an increasingly important issue. In particular, with the spread of market-linked electricity pricing schemes, it has become crucial to formulate economically efficient charging and discharging schedules that respond to market supply-demand conditions.
This study proposes an optimal operational method aimed at reducing costs and improving efficiency in a power grid system---comprising photovoltaic facilities, storage batteries, and self-consignment destinations---modeled on a factory located in Kochi Prefecture. Specifically, two mathematical models---a linear programming model and a dynamic programming model---are constructed to compare and verify their computational characteristics and practical applicability. The linear programming model derives exact solutions by setting multiple objective functions, such as minimizing electricity purchasing costs and minimizing net expenditure while accounting for revenue from surplus power sales. In contrast, the dynamic programming model discretizes battery capacity and is formulated through the Bellman equation based on the principle of optimality.
Experiments using real data from July 2025 show that although both models yield reasonable operational plans, distinct differences appear in their computational characteristics. The linear programming model is more sensitive to constraints, whereas the dynamic programming model demonstrates stable scalability, with computation time increasing linearly with the length of the planning period. In conclusion, this study establishes guidelines for model selections in the real environments: the linear programming model is effective for short-term planning and high-precision operations, while the dynamic programming model is preferable when long-term planning and predictable computation times are prioritized.


